\documentclass{amsart}
% For dealing with references we use the comment environment
\usepackage{verbatim}
\newenvironment{reference}{\comment}{\endcomment}
%\newenvironment{reference}{}{}
\newenvironment{slogan}{\comment}{\endcomment}
\newenvironment{history}{\comment}{\endcomment}

% For commutative diagrams we use Xy-pic
\usepackage[all]{xy}

% We use 2cell for 2-commutative diagrams.
\xyoption{2cell}
\UseAllTwocells

% We use multicol for the list of chapters between chapters
\usepackage{multicol}

% This is generally recommended for better output
\usepackage{lmodern}
\usepackage[T1]{fontenc}

% For cross-file-references

% Package for hypertext links:
\usepackage{hyperref}

% For any local file, say "hello.tex" you want to link to please

% Theorem environments.
%
\theoremstyle{plain}
\newtheorem{theorem}[subsection]{Theorem}
\newtheorem{proposition}[subsection]{Proposition}
\newtheorem{lemma}[subsection]{Lemma}

\theoremstyle{definition}
\newtheorem{definition}[subsection]{Definition}
\newtheorem{example}[subsection]{Example}
\newtheorem{exercise}[subsection]{Exercise}
\newtheorem{situation}[subsection]{Situation}

\theoremstyle{remark}
\newtheorem{remark}[subsection]{Remark}
\newtheorem{remarks}[subsection]{Remarks}

\numberwithin{equation}{subsection}

% Macros
%
\def\lim{\mathop{\mathrm{lim}}\nolimits}
\def\colim{\mathop{\mathrm{colim}}\nolimits}
\def\Spec{\mathop{\mathrm{Spec}}}
\def\Hom{\mathop{\mathrm{Hom}}\nolimits}
\def\Ext{\mathop{\mathrm{Ext}}\nolimits}
\def\SheafHom{\mathop{\mathcal{H}\!\mathit{om}}\nolimits}
\def\SheafExt{\mathop{\mathcal{E}\!\mathit{xt}}\nolimits}
\def\Sch{\mathit{Sch}}
\def\Mor{\mathop{\mathrm{Mor}}\nolimits}
\def\Ob{\mathop{\mathrm{Ob}}\nolimits}
\def\Sh{\mathop{\mathit{Sh}}\nolimits}
\def\NL{\mathop{N\!L}\nolimits}
\def\CH{\mathop{\mathrm{CH}}\nolimits}
\def\proetale{{pro\text{-}\acute{e}tale}}
\def\etale{{\acute{e}tale}}
\def\QCoh{\mathit{QCoh}}
\def\Ker{\mathop{\mathrm{Ker}}}
\def\Im{\mathop{\mathrm{Im}}}
\def\Coker{\mathop{\mathrm{Coker}}}
\def\Coim{\mathop{\mathrm{Coim}}}

% Boxtimes
%
\DeclareMathSymbol{\boxtimes}{\mathbin}{AMSa}{"02}

%
% Macros for moduli stacks/spaces
%
\def\QCohstack{\mathcal{QC}\!\mathit{oh}}
\def\Cohstack{\mathcal{C}\!\mathit{oh}}
\def\Spacesstack{\mathcal{S}\!\mathit{paces}}
\def\Quotfunctor{\mathrm{Quot}}
\def\Hilbfunctor{\mathrm{Hilb}}
\def\Curvesstack{\mathcal{C}\!\mathit{urves}}
\def\Polarizedstack{\mathcal{P}\!\mathit{olarized}}
\def\Complexesstack{\mathcal{C}\!\mathit{omplexes}}
% \Pic is the operator that assigns to X its picard group, usage \Pic(X)
% \Picardstack_{X/B} denotes the Picard stack of X over B
% \Picardfunctor_{X/B} denotes the Picard functor of X over B
\def\Pic{\mathop{\mathrm{Pic}}\nolimits}
\def\Picardstack{\mathcal{P}\!\mathit{ic}}
\def\Picardfunctor{\mathrm{Pic}}
\def\Deformationcategory{\mathcal{D}\!\mathit{ef}}

\setlength{\emergencystretch}{3em}
\begin{document}
\title[Stacks Project, Tag 0BFP]{578 --- Proper birational morphisms with reduced source are isomorphisms near normal target points admitting a quasi-finite source point}
\maketitle
\noindent Copyright (C) 2005--2025 Johan de Jong. Stacks Project authors. Source: \href{https://stacks.math.columbia.edu/tag/0BFP}{Tag 0BFP}.

\noindent Editorial difficulty note (not author text). Difficulty H1: The proof must turn the local quasi-finite point into a local isomorphism rather than merely identify a finite generic extension. A finite intermediate algebra inside the common function field is forced to equal the normal local ring, and a locally defined section plus properness and reducedness gives an isomorphism on a neighbourhood..

\noindent Editorial provenance note (not author text). History: excerpt from revision \texttt{a0444\allowbreak 6e57e\allowbreak c1fbc\allowbreak 252a8\allowbreak 71afc\allowbreak ec775\allowbreak 2fb28\allowbreak 07b14}, varieties.tex, lines 3299--3377. Reference presentation was adapted; original-author context and core support blocks were retained or added. The mathematical wording is unchanged. Licensed under GNU FDL 1.2 or later; no invariant sections or cover texts. See ../licenses/COPYING. Context blocks reproduce author definitions and supporting results; they are not additional counted problems.

\begin{definition}
\label{algebra-definition-quasi-finite}
Let $R \to S$ be a finite type ring map.
Let $\mathfrak q \subset S$ be a prime.
\begin{enumerate}
\item If the equivalent conditions of Lemma \href{https://stacks.math.columbia.edu/tag/00PK}{lemma isolated point fibre (Tag 00PK)}
are satisfied then we say $R \to S$ is {\it quasi-finite at $\mathfrak q$}.
\item We say a ring map $A \to B$ is {\it quasi-finite}
if it is of finite type and quasi-finite at all primes of $B$.
\end{enumerate}
\end{definition}

\begin{definition}
\label{morphisms-definition-quasi-finite}
\begin{reference}
\href{https://stacks.math.columbia.edu/bibliography}{Bibliography: EGA (II Definition 6.2.3)}
\end{reference}
Let $f : X \to S$ be a morphism of schemes.
\begin{enumerate}
\item We say that $f$ is {\it quasi-finite at a point $x \in X$}
if there exist an affine neighbourhood $\Spec(A) = U \subset X$
of $x$ and an affine open $\Spec(R) = V \subset S$ such that
$f(U) \subset V$, the ring map $R \to A$ is of finite type,
and $R \to A$ is quasi-finite at the prime of $A$ corresponding to $x$
(see above).
\item We say $f$ is {\it locally quasi-finite} if $f$ is
quasi-finite at every point $x$ of $X$.
\item We say that $f$ is {\it quasi-finite} if $f$ is of finite type
and every point $x$ is an isolated point of its fibre.
\end{enumerate}
\end{definition}

\begin{lemma}
\label{lemma-modification-normal-iso-over-codimension-1}
Let $X$ be a Noetherian scheme. Let $f : Y \to X$ be a birational proper
morphism of schemes with $Y$ reduced. Let $U \subset X$ be the
maximal open over which $f$ is an isomorphism. Then $U$ contains
\begin{enumerate}
\item every point of codimension $0$ in $X$,
\item every $x \in X$ of codimension $1$ on $X$ such that
$\mathcal{O}_{X, x}$ is a discrete valuation ring,
\item every $x \in X$ such that the fibre of $Y \to X$ over $x$ is
finite and such that $\mathcal{O}_{X, x}$ is normal, and
\item every $x \in X$ such that $f$ is quasi-finite at some
$y \in Y$ lying over $x$ and $\mathcal{O}_{X, x}$ is normal.
\end{enumerate}
\end{lemma}

\begin{proof}
Part (1) follows from Morphisms, Lemma
\href{https://stacks.math.columbia.edu/tag/0BAJ}{lemma birational isomorphism over dense open (Tag 0BAJ)}.
Part (2) follows from part (3) and Lemma \href{https://stacks.math.columbia.edu/tag/0AB7}{lemma finite in codim 1 (Tag 0AB7)}
(and the fact that finite morphisms have finite fibres).

\medskip\noindent
Part (3) follows from part (4) and
Morphisms, Lemma \href{https://stacks.math.columbia.edu/tag/02NG}{lemma finite fibre (Tag 02NG)}
but we will also give a direct proof.
Let $x \in X$ be as in (3). By
Cohomology of Schemes, Lemma
\href{https://stacks.math.columbia.edu/tag/02OH}{lemma proper finite fibre finite in neighbourhood (Tag 02OH)}
we may assume $f$ is finite. We may assume $X$ affine.
This reduces us to
the case of a finite birational morphism of Noetherian affine schemes
$Y \to X$ and $x \in X$ such that $\mathcal{O}_{X, x}$ is a
normal domain. Since $\mathcal{O}_{X, x}$ is a domain and $X$
is Noetherian, we may replace $X$ by an affine open of $x$ which
is integral. Then, since $Y \to X$ is birational and $Y$ is reduced
we see that $Y$ is integral. Writing $X = \Spec(A)$ and $Y = \Spec(B)$
we see that $A \subset B$ is a finite inclusion of domains having the same
field of fractions. If $\mathfrak p \subset A$ is the prime corresponding
to $x$, then $A_\mathfrak p$ being normal implies that
$A_\mathfrak p \subset B_\mathfrak p$ is an equality.
Since $B$ is a finite $A$-module, we see there exists an
$a \in A$, $a \not \in \mathfrak p$ such that $A_a \to B_a$
is an isomorphism.

\medskip\noindent
Let $x \in X$ and $y \in Y$ be as in (4). After replacing $X$
by an affine open neighbourhood we may assume $X = \Spec(A)$
and $A \subset \mathcal{O}_{X, x}$, see
Properties, Lemma \href{https://stacks.math.columbia.edu/tag/0BX3}{lemma ring affine open injective local ring (Tag 0BX3)}.
Then $A$ is a domain and hence $X$ is integral.
Since $f$ is birational and $Y$ is reduced
it follows that $Y$ is integral too. Consider the ring map
$\mathcal{O}_{X, x} \to \mathcal{O}_{Y, y}$. This is a ring map
which is essentially of finite type, the residue field extension
is finite, and $\dim(\mathcal{O}_{Y, y}/\mathfrak m_x\mathcal{O}_{Y, y}) = 0$
(to see this trace through the definitions of quasi-finite
maps in
Morphisms, Definition \ref{morphisms-definition-quasi-finite} and
Algebra, Definition \ref{algebra-definition-quasi-finite}). By
Algebra, Lemma \href{https://stacks.math.columbia.edu/tag/052V}{lemma essentially finite type fibre dim zero (Tag 052V)}
$\mathcal{O}_{Y, y}$ is the localization of a finite
$\mathcal{O}_{X, x}$-algebra $B$. Of course we may replace
$B$ by the image of $B$ in $\mathcal{O}_{Y, y}$ and assume
that $B$ is a domain with the same fraction field as $\mathcal{O}_{Y, y}$.
Then $\mathcal{O}_{X, x} \subset B$
have the same fraction field as $f$ is birational. Since
$\mathcal{O}_{X, x}$ is normal, we conclude that
$\mathcal{O}_{X, x} = B$ (because finite implies integral),
in particular, we see that $\mathcal{O}_{X, x} = \mathcal{O}_{Y, y}$. By
Morphisms, Lemma \href{https://stacks.math.columbia.edu/tag/0BX6}{lemma morphism defined local ring (Tag 0BX6)}
after shrinking $X$ we may assume there is a section
$X \to Y$ of $f$ mapping $x$ to $y$ and inducing the given
isomorphism on local rings. Since $X \to Y$ is closed
(by Schemes, Lemma \href{https://stacks.math.columbia.edu/tag/01KT}{lemma section immersion (Tag 01KT)})
necessarily maps the generic point of $X$ to the generic point of $Y$
it follows that the image of $X \to Y$ is $Y$.
Then $Y = X$ and we've proved what we wanted to show.
\end{proof}
\end{document}
